\chapter{Version Control and Git's Model}\label{ch:git-model}

Nearly all confusion about Git is confusion about its model, not its commands.
People who have memorised \cmd{git pull} and \cmd{git push} without ever being
told what a commit \emph{is} will eventually hit a merge conflict, a detached
\texttt{HEAD} or a bad rebase, and have no way to reason about what they are
looking at.

So this chapter runs no useful commands at all. It is about four ideas ---
snapshots, objects, the three trees, and refs --- and everything in
Part~\ref{part:git} follows from them.

\section{Notation for Part V}\label{sec:git-notation}

Git is driven from the shell, so transcripts keep the convention of
\autoref{sec:sh-notation}: \texttt{\$} is what you type, everything else is
output.

\begin{shell}
$ git status
On branch main
nothing to commit, working tree clean
\end{shell}

Commit graphs are drawn with the newest commit on the right, which is the
opposite of how \cmd{git log} prints them and the same as how everyone draws
them on a whiteboard:

\begin{center}
	\begin{tikzpicture}[
			x=1mm, y=1mm, font=\footnotesize\ttfamily,
			c/.style={draw=Gray!60, circle, inner sep=2pt, fill=Gray!5, minimum size=7mm},
			ref/.style={draw=SeaGreen!60!black, rounded corners=2pt, inner sep=2.5pt,
					fill=SeaGreen!8, font=\scriptsize\ttfamily},
			e/.style={-{Stealth[length=4pt]}, Gray!70!black}
		]
		\node[c] (a) at (  0, 0) {A};
		\node[c] (b) at ( 22, 0) {B};
		\node[c] (d) at ( 44, 0) {C};
		\draw[e] (b) -- (a);
		\draw[e] (d) -- (b);
		\node[ref] (m) at ( 44, 16) {main};
		\node[ref] (h) at ( 44, 28) {HEAD};
		\draw[e, SeaGreen!60!black] (m) -- (d);
		\draw[e, SeaGreen!60!black] (h) -- (m);
	\end{tikzpicture}
\end{center}

\begin{keys}[0.26]
	\texttt{A}, \texttt{B}, \texttt{C} & Commits, oldest on the left            \\
	\cmd{main}, \cmd{feature}          & Branches --- names pointing at a commit \\
	\cmd{HEAD}                         & Where you are now                       \\
	\texttt{a1b2c3d}                   & An abbreviated commit hash              \\
	\texttt{<sha>}, \texttt{<branch>}  & Placeholders, as in \autoref{sec:sh-notation} \\
\end{keys}

\begin{note}
	The arrows point from a commit to its \emph{parent}, backwards in time. That is
	the direction the data actually goes: a commit records what came before it, and
	nothing records what comes after. Almost every surprising property of Git ---
	why history is cheap to rewrite, why a branch is free, why a lost commit is
	still there --- follows from that one arrow direction.
\end{note}

\section{Why Version Control}\label{sec:git-why}

The obvious answer is backups, and it is the least interesting one. What a
version control system actually gives you:

\begin{keys}[0.24]
	\textbf{A safety net}    & Any change can be undone, including the one you made three weeks ago \\
	\textbf{A record}        & Every line has an author, a date, and a message saying why           \\
	\textbf{Parallel work}   & Several people, or several of your own ideas, on the same files      \\
	\textbf{Bisection}       & ``It worked in March'' becomes a mechanical search (\autoref{sec:git-bisect}) \\
	\textbf{A shared truth}  & One authoritative history, rather than \file{report\_final\_v3\_FINAL.tex} \\
\end{keys}

\begin{intuition}
	The real change is psychological. With history you can afford to try things ---
	delete a section, restructure a module, rewrite an algorithm --- because the
	cost of being wrong is one command. Without it, every experiment is a
	negotiation with your own nerve.
\end{intuition}

\section{Snapshots, Not Diffs}\label{sec:git-snapshots}

This is the design decision that makes Git different from what came before.

\begin{definition}[Commit]
	A complete \emph{snapshot} of every tracked file at one moment, together with
	its author, its date, its message, and a pointer to the commit or commits that
	came before it.
\end{definition}

Older systems --- CVS, Subversion --- stored a list of files and, for each, the
sequence of changes made to it. Git stores the whole tree, every time.

\begin{intuition}
	That sounds wasteful and is not: a file that did not change between two commits
	is not stored twice, it is \emph{referenced} twice, because Git identifies
	content by its hash. Storing snapshots rather than diffs means any commit can
	be reconstructed without replaying history, which is why \cmd{git checkout} of
	a five-year-old commit is instant.
\end{intuition}

\begin{note}
	\cmd{git diff} shows you a diff, so it is easy to conclude that Git stores
	diffs. It does not: the diff is computed on demand by comparing two snapshots.
	This is why Git can show a diff between \emph{any} two commits, related or not.
\end{note}

\section{The Object Database}\label{sec:git-objects}

Everything Git stores is one of four kinds of object, each named by the SHA-1
hash of its own contents.

\begin{keys}[0.20]
	\vocab{blob}   & The contents of a file. No name, no permissions --- just bytes \\
	\vocab{tree}   & A directory: a list of names, each with a mode and a pointer to a blob or another tree \\
	\vocab{commit} & A pointer to one tree, plus parents, author, date and message  \\
	\vocab{tag}    & An annotated name for a commit                                 \\
\end{keys}

\begin{center}
	\begin{tikzpicture}[
			x=1mm, y=1mm, font=\footnotesize\ttfamily,
			o/.style={draw=Gray!60, rounded corners=2pt, inner sep=3.5pt, fill=Gray!5},
			e/.style={-{Stealth[length=4pt]}, Gray!70!black}
		]
		\node[o] (c)  at (  0,  0) {commit};
		\node[o] (t)  at ( 30,  0) {tree};
		\node[o] (b1) at ( 66, 10) {blob \textit{main.c}};
		\node[o] (t2) at ( 66, -8) {tree \textit{src/}};
		\node[o] (b2) at (104, -8) {blob \textit{util.c}};
		\draw[e] (c) -- (t);
		\draw[e] (t) -- (b1);
		\draw[e] (t) -- (t2);
		\draw[e] (t2) -- (b2);
	\end{tikzpicture}
\end{center}

\begin{eg}
	You can look at any of them directly:
\begin{shell}
$ git cat-file -t HEAD           # what kind of object is this?
commit
$ git cat-file -p HEAD           # print it
tree 4b825dc642cb6eb9a060e54bf8d69288fbee4904
parent 8f3a1c9e2b7d4a5f6c8e9d0a1b2c3d4e5f6a7b8c
author Kon Yi <you@example.com> 1757808000 +0800
committer Kon Yi <you@example.com> 1757808000 +0800

    Week1 Finish
$ git cat-file -p HEAD^{tree}    # the directory it points at
100644 blob a3f2e1d...    README.md
040000 tree 9c4b8a2...    Chapters
\end{shell}
	That is the entire format. A commit is a small text file naming a tree; a tree
	is a list naming blobs and other trees.
\end{eg}

\begin{definition}[Content addressing]
	An object's name \emph{is} the hash of its contents. Two identical files
	anywhere in history are one blob. Changing a byte changes the hash, which
	changes the tree that names it, which changes the commit --- so a commit hash
	certifies the entire tree beneath it.
\end{definition}

\begin{moral}
	This is why history cannot be quietly altered: editing an old commit gives it a
	new hash, which changes every descendant's hash too. ``Rewriting history''
	(\autoref{sec:git-rewrite}) does not modify commits --- it makes new ones and
	moves the branch.
\end{moral}

\section{The Three Trees}\label{sec:git-threetrees}

Git manages three versions of your project at once, and almost every confusing
Git message is about the boundary between two of them.

\begin{center}
	\begin{tikzpicture}[
			x=1mm, y=1mm, font=\footnotesize\sffamily,
			t/.style={draw=Gray!60, rounded corners=3pt, inner sep=5pt,
					fill=Gray!5, minimum width=30mm, align=center},
			e/.style={-{Stealth[length=4pt]}, Gray!70!black},
			lbl/.style={font=\scriptsize\ttfamily, Gray!60!black}
		]
		\node[t] (w) at (  0, 0) {Working tree\\[1pt]\scriptsize the files on disk};
		\node[t] (i) at ( 52, 0) {Index\\[1pt]\scriptsize the next commit};
		\node[t] (h) at (104, 0) {HEAD\\[1pt]\scriptsize the last commit};
		\draw[e] (w) -- node[lbl, above]{git add} (i);
		\draw[e] (i) -- node[lbl, above]{git commit} (h);
		\draw[e, SeaGreen!60!black] (h.south) to[out=-90,in=-90,looseness=0.45]
		node[lbl, below, pos=0.5]{git checkout / git restore} (w.south);
	\end{tikzpicture}
\end{center}

\begin{keys}[0.24]
	\vocab{Working tree} & The actual files in your directory. What your editor sees \\
	\vocab{Index} (or \vocab{staging area}) & A draft of the next commit. \cmd{git add} copies from the working tree into it \\
	\vocab{HEAD}         & The commit you last committed or checked out             \\
\end{keys}

\begin{intuition}
	The index is the feature people find pointless until they need it, and then
	cannot do without. It lets you decide \emph{what} to commit separately from
	\emph{what you have changed} --- so a session that fixed a bug and also
	reformatted a function becomes two clean commits instead of one muddled one.
	\cmd{git add -p} (\autoref{sec:git-add}) is where this pays off.
\end{intuition}

\begin{note}
	Now \cmd{git status} reads as a report on two boundaries:
	\emph{``Changes to be committed''} is index versus \texttt{HEAD};
	\emph{``Changes not staged for commit''} is working tree versus index; and
	\emph{``Untracked files''} are in neither. \cmd{git diff} compares working tree
	with index, and \cmd{git diff --staged} compares index with \texttt{HEAD}.
	Nothing else needs to be memorised --- it all falls out of the diagram.
\end{note}

\section{Refs, \texorpdfstring{\texttt{HEAD}}{HEAD}, and What a Branch Really Is}\label{sec:git-refs}

\begin{definition}[Branch]
	A file containing forty hex digits. That is not a simplification: a branch is a
	\vocab{ref} --- a name pointing at one commit --- stored in
	\file{.git/refs/heads/}, and committing moves the pointer forward.
\end{definition}

\begin{shell}
$ cat .git/refs/heads/main
bc91a17e5f3d2c1b0a9876543210fedcba987654
$ cat .git/HEAD
ref: refs/heads/main
\end{shell}

\begin{moral}
	Branches are free. Creating one writes 41 bytes; deleting one removes them.
	Nothing is copied, no files move, and there is no reason to hesitate before
	making one. If you come from Subversion, where a branch was a server-side copy
	of the tree, this is the single biggest adjustment.
\end{moral}

\texttt{HEAD} is one more level of indirection: it normally holds the \emph{name}
of a branch rather than a commit. So committing updates \texttt{HEAD}'s branch,
and \cmd{git switch} rewrites \texttt{HEAD} to name a different one.

\begin{definition}[Detached \texttt{HEAD}]
	The state where \texttt{HEAD} holds a commit hash directly instead of a branch
	name --- after \cmd{git checkout <sha>}, or checking out a tag. Commits made
	there belong to no branch, so moving away leaves them unreferenced and
	apparently lost.
\end{definition}

\begin{note}
	This is the state that alarms people, and it is harmless as long as you
	understand it. Nothing is broken; you are simply standing on a commit with no
	name attached. \cmd{git switch -c fix} gives the current position a name and
	you are back on ordinary ground --- and even if you walk away first,
	\autoref{sec:git-reflog} finds it again.
\end{note}

\begin{keys}[0.22]
	\cmd{HEAD}          & Where you are                                       \\
	\cmd{HEAD\textasciitilde{}1}, \cmd{HEAD\textasciitilde{}3} & One, three commits back along first parents \\
	\cmd{HEAD\textasciicircum{}}, \cmd{HEAD\textasciicircum{}2} & The first, second parent of a merge \\
	\cmd{main@\{yesterday\}} & Where \cmd{main} pointed yesterday               \\
	\cmd{a1b2c3d}       & A commit by hash; seven characters is usually enough \\
	\cmd{v1.0}          & A tag                                                \\
	\cmd{origin/main}   & A \vocab{remote-tracking} branch: where \cmd{main} was on the remote when you last fetched \\
\end{keys}

\section{First-Time Setup}\label{sec:git-config}

\begin{shell}
$ git config --global user.name  "Kon Yi"
$ git config --global user.email "you@example.com"
$ git config --global init.defaultBranch main
$ git config --global core.editor "nvim"
$ git config --global pull.rebase true         # see sec. 33.4
$ git config --global --list                   # check what you have
\end{shell}

\begin{keys}[0.30]
	\cmd{-{}-global}        & \file{\textasciitilde/.gitconfig} --- all your repositories \\
	\cmd{-{}-local}         & \file{.git/config} --- this repository only. The default    \\
	\cmd{-{}-system}        & Every user on the machine                                  \\
	\cmd{-{}-list -{}-show-origin} & Every setting, and which file set it                \\
\end{keys}

\begin{note}
	The name and email are baked into every commit you make and cannot be changed
	afterwards without rewriting history. Set them before your first commit. If
	you use a different address for work, a conditional include keeps them apart:
\begin{conf}
# ~/.gitconfig
[user]
    name = Kon Yi
    email = personal@example.com

[includeIf "gitdir:~/work/"]
    path = ~/.gitconfig-work
\end{conf}
\end{note}

\begin{note}[This machine]
	Git 2.43 is installed, with \cmd{user.name} set to \texttt{Konyi0613} and
	credentials handled by the GitHub CLI --- \cmd{gh auth git-credential} is
	registered as the credential helper for \texttt{github.com}, so pushing needs
	no password prompt. \autoref{sec:gh-auth} covers what that arrangement is
	doing.
\end{note}

\begin{tryit}
	In any repository you have, run \cmd{git cat-file -p HEAD}, then
	\cmd{git cat-file -p HEAD\textasciicircum{}\{tree\}}, then \cmd{cat .git/HEAD}.
	Three commands, and the abstraction is gone: a commit really is a short text
	file naming a tree, and a branch really is a file with a hash in it.
\end{tryit}
